\newpage
\pagestyle{fancy}
\fancyhf{}
\fancyhead[R]{\thepage}
\chapter{PENDAHULUAN} \label{Bab I}

\section{Latar Belakang} \label{I.latar belakang}
Lanjut usia (lansia) merupakan fase alami dalam siklus kehidupan manusia yang akan dialami oleh setiap individu. \textit{World Health Organization} (WHO) mengelompokan populasi yang memasuki usia 60 tahun  keatas sebagai lansia \cite{Gusti2023}. Populasi tersebut memiliki kecenderungan penurunan fungsi karakteristik fisik dan kondisi kesehatan  yang semakin signifikan seiring bertambahnya usia. Perubahan fisiologis pada lansia seperti berkurangnya massa otot, penurunan kepadatan tulang, dan gangguan keseimbangan membuat mereka lebih rentan terhadap berbagai masalah kesehatan \cite{Colon2018}. Kondisi ini membuat lansia lebih rentan terhadap masalah kesehatan dan keselamatan, seperti risiko terjatuh. \par

Jatuh merupakan suatu kondisi di mana seseorang tanpa sadar tiba- tiba berbaring atau terduduk ke permukaan apapun, seperti tanah, lantai, tempat tidur, atau permukaan lainnya. Jatuh juga didefinisikan sebagai suatu kejadian dimana seseorang tiba-tiba jatuh baik disengaja maupun tidak disengaja yang mengakibatkan luka atau cedera pada orang tersebut \cite{Forrest2012}. Orang dewasa dan lansia sangat rentan mengalami cedera akibat jatuh karena berbagai faktor fisik dan lingkungan. Penurunan fungsi sensorik, gangguan keseimbangan, berkurangnya massa otot, dan efek samping obat-obatan menjadi penyebab utama tingginya risiko jatuh pada populasi ini \cite{Colon2024}.\par

Menurut data statistik dari WHO, sekitar 37,3 juta kejadian jatuh yang memerlukan perhatian medis, bahkan terdapat 684.000 orang meninggal akibat terjatuh di seluruh dunia, di mana lebih dari 80\% di antaranya terjadi di negara-negara berkembang \cite{Who2021}. Studi yang dilakukan \textit{Australian Centre for Evidence Based Care} (ACEBAC) menyatakan bahwa Lansia yang tinggal di panti jompo lebih sering jatuh dari pada lansia yang tinggal dirumah mencapai 30-50\% setiap tahun dan meningkat 40\% mengalami jatuh berulang \cite{ACEBAC2023}. Berdasarkan data survei \textit{Indonesian Family Life Survey} (IFLS), Angka kejadian jatuh pada lansia di Indonesia yang berusia 65 tahun ke atas sekitar 30\% dan usia 80 tahun ke atas sekitar 50\% \cite{Gusti2023}.\par

Di panti jompo atau lingkungan perawatan lansia, terdapat periode waktu di mana subjek lansia sedang sendirian dan tidak dalam pengawasan langsung oleh pengasuh atau perawat. Sebuah studi yang dilakukan Albasha, menunjukkan bahwa tingkat jatuh di fasilitas perawatan jangka panjang, termasuk panti jompo, jauh lebih tinggi dibandingkan dengan lingkungan pribadi atau keluarga, dengan rata-rata 1,7 atau hampir 2 kali kejadian jatuh per penghuni setiap tahun. Jatuh lebih sering terjadi pada malam hari atau selama pergantian \textit{shift} karena rendahnya tingkat staf yang bertugas, mengakibatkan kurangnya pengawasan langsung terhadap penghuni \cite{Albasha2024}. Kondisi ini meningkatkan risiko keterlambatan penanganan karena ketiadaan sistem peringatan dini yang efektif untuk mendeteksi kejadian jatuh.\par

Keterlambatan penanganan pada lansia yang mengalami jatuh dapat menyebabkan komplikasi serius hingga kematian. Menurut penelitian oleh Fleming, fungsi otot lansia yang terjatuh akan menurun sehingga mengurangi kemampuan untuk bangun. Jika tidak mampu bangun selama lebih dari satu jam akan mengalami dehidrasi, hipotermia, pneumonia, dan ulkus tekanan. dilakukan eksperiman dengan 20 orang yang berbaring di lantai selama lebih dari satu jam setelah jatuh, 36\% (5 dari 14) pindah ke perawatan jangka panjang dalam satu tahun, dan 53\% (8 dari 15) pindah ke perawatan jangka panjang pada akhir masa penelitian \cite{Fleming2008}. Hal ini menggarisbawahi dibutuhkannya sistem pendeteksi jatuh yang dapat berfungsi 24 jam sehari, terutama di area-area yang tidak selalu dalam pengawasan langsung.\par

Permasalahannya secara umum adalah bagaimana cara memberikan peringatan langsung kepada pengasuh atau petugas medis ketika subjek lansia mengalami jatuh. Sejauh ini teknologi deteksi jatuh telah dikembangkan dengan 3 pendekatan, antara lain \textit{wearable sensors}, \textit{vision atau computer visison} dan gabungan keduanya. Pada pendekatan \textit{wearable}, seperti penggunaan perangkat sensor \textit{accelerometer} dan IMU (\textit{Inertial Measurement Unit}) memungkinkan pengumpulan data yang akurat selama aktivitas sehari-hari termasuk kejadian jatuh. Namun terdapat beberapa kelemahan yang perlu diperhatikan seperti sensor yang dipasang pada tubuh, seperti di pergelangan tangan atau pinggang, harus nyaman untuk dipakai sehari-hari, atau pengguna akan enggan menggunakannya. Selain itu terdapat ketergantungan pada kalibrasi membutuhkan penyesuaian manual, serta perangkat \textit{wearable} memiliki masa pakai baterai yang terbatas sehingga ada kemungkinan di mana kejadian jatuh terjadi pada saat daya baterai habis \cite{Mubashir2013}.

Berbagai kekurangan dari pendekatan \textit{wearable} menunjukkan kesempatan pada pendekatan \textit{computer vision} atau visi komputer, yaitu bidang teknologi yang berfokus pada pengembangan algoritma dan sistem yang memungkinkan komputer untuk memperoleh, memproses, dan memahami informasi dari gambar atau video secara otomatis \cite{Negoro2023}. Teknologi ini meniru fungsi penglihatan manusia dengan cara mengekstrak informasi visual dari input yang tersedia dan mengubahnya menjadi format yang dapat diproses lebih lanjut oleh sistem komputer. Demi terwujudnya prinsip ini dibutuhkan teknik pengembangan algoritma seperti \textit{Deep Learning} yang dapat menjadi inti dari sistem deteksi jatuh.

\textit{Deep learning} adalah metode pembelajaran yang mengandalkan jaringan saraf tiruan dalam menganalisis dan memprediksi data melalui banyak lapisan pemrosesan \cite{LeCun2015}. \textit{Deep learning} memanfaatkan arsitektur jaringan yang kompleks untuk mengekstraksi fitur dari data, yang dapat mengarah pada peningkatan kinerja dalam berbagai aplikasi, seperti pengenalan suara, pengenalan objek visual, dan berbagai aplikasi lain di bidang ilmu kehidupan dan industri. Beberapa sistem deteksi jatuh berbasis \textit{Deep Learning} telah dikembangkan sebelumnya, namun masih memiliki berbagai keterbatasan. Salah satu penelitian terkini yang relevan dalam bidang ini yaitu pengembangkan model deteksi jatuh berbasis visi dengan memanfaatkan \textit{Convolutional Neural Network} (CNN) dan \textit{Long Short-Term Memory} (LSTM) untuk mengelola informasi temporal dari urutan gambar \cite{Chhetri2021}. Meskipun inovatif, penelitian ini juga menunjukkan beberapa kelemahan signifikan yang perlu dicermati. 

Salah satu kelemahan utama sistem tersebut terletak pada kebutuhan akan daya komputasi yang tinggi untuk menjalankan model deep learning untuk pengolahan citra. Menurut Chhetri, implementasi teknik \textit{deep learning} pada model deteksi jatuh berbasis \textit{computer vision} memerlukan sumber daya komputasi yang besar, yang mungkin tidak selalu tersedia dalam perangkat keras sederhana atau untuk penggunaan dalam skala besar \cite{Chhetri2021}. Ini bisa menjadi kendala signifikan ketika mencoba menerapkan sistem pada perangkat mobile yang memiliki keterbatasan energi dan proses. \par

Pengolahan citra menggunakan MediaPipe \textit{pose landmark} menawarkan solusi yang lebih efisien untuk deteksi jatuh. MediaPipe adalah framework \textit{open-source} yang dikembangkan oleh Google untuk analisis pose tubuh atau wajah manusia melalui penerapan deteksi \textit{landmark}, yang menggunakan model \textit{deep learning} untuk mengekstraksi titik-titik kunci (\textit{keypoints}) pada tubuh manusia dari gambar atau video. Dibandingkan dengan pendekatan pengolahan citra konvensional, model pose landmark menghasilkan data graph dengan ukuran yang jauh lebih kecil karena hanya menyimpan koordinat titik-titik kunci, bukan seluruh citra \cite{Lugaresi2019}. Penelitian terbaru di tahun 2024 menunjukkan bahwa optimasi arsitektur MediaPipe mampu mengurangi waktu pemrosesan per frame hingga 30\%, sambil meningkatkan akurasi estimasi pose sebesar 20\% berdasarkan metrik \textit{Intersection over Union} (IoU) \cite{Sengar2024}. MediaPipe menunjukkan peningkatan ketahanan terhadap variasi pencahayaan dan oklusi parsial dengan memprioritaskan deteksi \textit{keypoint} di atas analisis tingkat piksel, memanfaatkan jaringan resolusi tinggi (HRNet) untuk mempertahankan akurasi bahkan ketika bagian tubuh dikaburkan \cite{Raj2024}.\par

\textit{Graph Convolutional Network} (GCN) adalah metode \textit{deep learning} yang dapat digunakan untuk memproses data berbentuk graf, seperti \textit{pose landmark} tubuh manusia. GCN pertama kali diperkenalkan oleh Kipf dan Welling sebagai metode untuk mengklasifikasikan \textit{node} dalam graf besar dengan memanfaatkan informasi topologi graf dan fitur \textit{node} \cite{Kipf2016}. Dalam konteks deteksi jatuh, Kerangka manusia dapat direpresentasikan sebagai sebuah graf dengan sendi-sendi sebagai \textit{node} dan tulang-tulang sebagai \textit{edge}, lalu membuat pemodelan GCN untuk mempelajari hubungan spasial dan temporal antar titik-titik kunci tubuh manusia, sehingga dapat mengenali pola postur tubuh yang mengindikasikan posis jatuh \cite{Zahan2025}. Penelitian yang dilakukan Zhang et al. (2023) menunjukkan bahwa pendekatan \textit{Graph Convolutional Network} berbasis \textit{pose landmark} mencapai akurasi 96,3\% pada dataset MCF, lebih tinggi dibandingkan metode konvensional yang hanya mencapai 93,4\% \cite{Zhang2023}.\par

Penelitian ini berfokus pada pengembangan model deteksi jatuh dengan menerapkan \textit{Graph Convolutional Network} (GCN) berbasis \textit{pose landmark}. Pengembangan model ini diharapkan mampu mengidentifikasi posisi jatuh yang lebih akurat sehingga menjadi dasar pengembangan sistem yang baru untuk peringatan dini kepada pengasuh atau petugas medis ketika subjek lansia mengalami jatuh, sehingga pertolongan dapat segera diberikan. Penelitian ini juga diharapkan dapat menjadi dasar pengembangan teknologi yang lebih canggih sehingga dapat mengurangi risiko komplikasi serius dan kematian akibat keterlambatan penanganan pasca-jatuh pada lansia di panti jompo, rumah sakit, atau bahkan rumah pribadi di masa depan, seiring dengan semakin meningkatnya populasi lansia di Indonesia dan dunia.\par

\section{Rumusan masalah} \label{I.Rumusan Masalah}
Berdasarkan latar belakang yang telah dijelaskan, maka rumusan masalah dalam penelitian ini sebagai berikut.
\begin{enumerate}[noitemsep]
    \item Bagaimana cara mengembangkan model menggunakan \textit{Graph Convolutional Network} (GCN) berbasis \textit{pose landmark} Mediapipe untuk mendeteksi posisi jatuh pada lansia?
    \item Bagaimana hasil evaluasi performa, terutama berdasarkan akurasi dari \textit{confusion matrix} pada model deteksi posisi jatuh GCN berbasis \textit{pose landmark} Mediapipe dengan membandingkan konfigurasi \textit{hyperparameter} default dan hasil studi ablasi?
\end{enumerate}


\section{Tujuan} \label{I.Tujuan}
Tujuan dari penelitian ini sebagai berikut.
\begin{enumerate}[noitemsep]
    \item Mengembangkan model deteksi posisi jatuh pada lansia menggunakan GCN berbasis pose landmark Mediapipe.
    \item Melakukan evaluasi performa, terutama berdasarkan akurasi dari \textit{confusion matrix} pada model deteksi posisi jatuh GCN berbasis \textit{pose landmark} Mediapipe dengan membandingkan Hyperparameter default dan studi ablasi.
\end{enumerate}

\section{Batasan Masalah} \label{I.Batasan}
Agar penelitian ini lebih terfokus, beberapa batasan masalah yang ditetapkan sebagai berikut
\begin{enumerate}[noitemsep]
    \item Penelitian ini hanya menggunakan dataset simulasi jatuh yang telah dibuat oleh Baldewijns (2016), terdiri dari 275 video jatuh dan 85 video tidak jatuh.
    \item Penelitian ini hanya berfokus pada pipeline pembangunan model dan eksplorasi hyperparameter  saja, tidak melakukan perbandingan efisiensi komputasi dengan penilitian terdahulu.
    \item Frame citra yang diproses adalah citra yang berisi 1 orang dengan pose keseluruhan tubuh yang jelas dan tidak tertutup oleh objek apapun.
    \item Penelitian ini tidak menerapkan multimodal dalam pengembangan dan pelatihan model.
    \item Hanya berfokus pada pengolahan dataset kumpulan video dimana tiap video akan diekstrak menjadi frame citra, lalu penerapan \textit{framework} Mediapipe \textit{Pose Landmark} untuk mendeteksi pose manusia dalam tiap citra dan terakhir pembangunan model berbasis GCN untuk memprediksi posisi jatuh dan tidak jatuh pada manusia.
    \item Framework Mediapipe Pose Landmark hanya sebatas pemakaian produk.
    \item Model yang dikembangkan hanya melakukan klasifikasi antara 2 kelas saja, jatuh dan tidak jatuh.
\end{enumerate}

\section{Manfaat Penelitian} \label{I.Manfaat}
Penelitian ini diharapkan memberikan manfaat sebagai berikut:
\begin{enumerate}[noitemsep]
    \item Memberikan informasi bagaimana proses pengembangan model deteksi kejatuhan pada lansia menggunakan Graph Convolutional Network (GCN) berbasis pose landmark Mediapipe.
    \item Mempelajari dan mengevaluasi performa model Graph Convolutional Network (GCN) berbasis pose landmark Mediapipe agar menjadi bahan referensi untuk pengembangan model deteksi posisi jatuh pada lansia di masa depan.
\end{enumerate}

\section{Sistematika Penulisan} \label{I.Sistematika}
Pada subbab ini akan menjelaskan isi dari setiap bab dalam laporan penelitian ini secara umum.  Adapun sistematika penulisan dalam laporan ini adalah sebagai berikut.\par
\subsection{Bab I Pendahuluan}
Bab ini memaparkan latar belakang, rumusan masalah, tujuan penelitian, batasan masalah, manfaat penelitian, dan sistematika penulisan. Permasalahan penelitian dirumuskan untuk memberikan fokus penelitian yang jelas, sementara tujuan penelitian menentukan pencapaian yang diharapkan. Batasan masalah mengarahkan ruang lingkup penelitian, dan manfaat penelitian diharapkan memberikan kontribusi positif baik teori maupun praktik.\par

\subsection{Bab II Tinjauan Pustaka}
Bab ini mengulas landasan teori, tinjauan pustaka, serta defenisi teori dan teknologi yang digunakan dalam penelitian, yang mencakup konsep-konsep dasar yang relevan untuk memahami topik yang dibahas. Dalam bagian ini, akan dijelaskan mengenai pose landmark dan Graph Convolutional Network (GCN) yang menjadi teknologi utama yang diterapkan dalam penelitian ini. Penjelasan mengenai pose landmark akan menguraikan cara teknologi ini mendeteksi dan menganalisis posisi tubuh manusia, sementara GCN akan dibahas dalam konteks bagaimana jaringan konvolusional graf dapat digunakan untuk memproses data spasial yang kompleks. Selain itu, studi literatur terkait deteksi kejatuhan akan dibahas untuk memberikan gambaran mengenai penelitian terdahulu yang relevan dan bagaimana temuan-temuan tersebut mendasari pengembangan penelitian ini.\par

\subsection{Bab III Metode Penelitian}
Bab ini menjelaskan secara rinci metode penelitian yang digunakan, yang mencakup langkah-langkah sistematis dalam proses penelitian. Pertama, dijelaskan mengenai pengumpulan data yang dilakukan untuk memperoleh informasi yang diperlukan, dengan menggambarkan teknik pengumpulan data serta sumber-sumber yang relevan untuk mendukung penelitian ini. Selanjutnya, proses prapemrosesan data akan diuraikan agar siap digunakan dalam analisis lebih lanjut. Implementasi metode GCN juga akan dijelaskan secara mendalam, mencakup bagaimana GCN diterapkan untuk memproses data yang diperoleh dan mengidentifikasi pola-pola dalam deteksi posisi jatuh, serta pengoptimalan model untuk mencapai hasil yang akurat dan efisien.\par

\subsection{Bab IV Hasil dan Pembahasan}
Bab ini berisi hasil penelitian yang diperoleh melalui penerapan metode yang telah dijelaskan sebelumnya, serta analisis mendalam terhadap temuan-temuan yang didapatkan. Dalam bagian ini, akan dipaparkan hasil eksperimen yang dilakukan untuk menguji kemampuan model dalam mendeteksi kejadian jatuh dan tidak jatuh, termasuk visualisasi hasil deteksi yang dihasilkan oleh model. Evaluasi performa model akan dilakukan dengan menggunakan metrik yang relevan, seperti akurasi, presisi, \textit{recall}, dan \textit{F1-score}, untuk menilai sejauh mana model dapat membedakan antara kedua kondisi tersebut. \par

\subsection{Bab V Kesimpulan dan Saran}
Bab ini menyajikan kesimpulan yang merangkum temuan-temuan utama dari penelitian yang telah dilakukan, serta implikasi dari hasil yang diperoleh dalam pengembangan deteksi posisi jatuh menggunakan metode GCN berbasis \textit{pose landmark}. Kesimpulan ini mengidentifikasi pencapaian utama penelitian, termasuk keberhasilan model dalam mendeteksi kejadian jatuh dan tidak jatuh, serta evaluasi terhadap efektivitas metode yang diterapkan. Selain itu, bab ini juga menyajikan saran untuk pengembangan lebih lanjut, meliputi potensi peningkatan model dengan memanfaatkan data tambahan atau teknik lain yang dapat meningkatkan akurasi deteksi. Saran ini juga mencakup arahan untuk penelitian lanjutan, serta aplikasi teknologi yang lebih luas dalam bidang deteksi kecelakaan pada lansia atau di lingkungan lain yang relevan, guna memberikan manfaat yang lebih besar bagi masyarakat.\par
